\documentclass[11pt]{article}

\usepackage[a4paper,margin=0.85in]{geometry}
\usepackage{amsmath,amssymb,amsthm,mathtools}
\usepackage[hidelinks]{hyperref}

\newtheorem{theorem}{Theorem}
\newtheorem{corollary}[theorem]{Corollary}
\newtheorem{remark}[theorem]{Remark}
\newcommand{\R}{\mathbb R}

\title{Critical exponential-moment dissipation and a no-go theorem\\
for repeating packets in nonlocal Fisher--KPP fronts}
\author{Hainan Zhao}
\date{Research note, August 2026}

\begin{document}
\maketitle

\begin{abstract}
For the one-dimensional nonlocal Fisher--KPP equation
\[
 u_t=D u_{xx}+u(1-K*u),
\]
with an even nonnegative competition kernel, we derive an exact dissipation
identity for every exponentially weighted mass after removing its linear
growth.  At the critical Fisher--KPP tilt, this gives a strict Lyapunov
quantity and an exact formula for the phase delay caused by nonlocal
competition.  A nonzero state with finite critical exponential moment that
repeats after a spatial translation must travel strictly slower than the
pulled speed $2\sqrt D$.  Thus a periodic-shedding description at pulled speed
cannot identify the complete localized state from one cycle to the next.  A
critical leading tail must be separated from the deposited wake, or an
additional competition-memory variable must be retained.  The result does
not select a wake wavelength or prove the proposed large-time pattern behind
the front.
\end{abstract}

\section{Setting and claim boundary}

Let $D>0$ and let $K\in L^1(\R)$ be even, nonnegative, nonzero, and normalized
by $\int_{\R}K=1$.  Consider
\begin{equation}
 u_t=D u_{xx}+u(1-K*u)
\end{equation}
on $\R$.  We assume that $u$ is a nonzero nonnegative classical solution on
$[t_0,t_1]$ and that the weighted integrals and integrations by parts below
are justified.  Bounded compactly supported initial data satisfy the needed
conditions at positive times under standard well-posedness hypotheses.  For
the top-hat model below, the required global bounded classical solution is
supplied by Hamel and Ryzhik~\cite[Theorem~1.2]{HamelRyzhik2014}.

The motivating special case is the radius-one top-hat kernel
\begin{equation}
 K(x)=\frac12\mathbf 1_{[-1,1]}(x),
\end{equation}
for which recent asymptotic and numerical work displays patterned wakes
behind an invasion front~\cite{NeedhamEtAl2025}.  Our conclusion concerns
exactly repeating, exponentially localized packets.  It does not exclude a
pulled front whose critical weighted moment diverges, nor a stationary wake
behind a separate leading edge.

\section{An exact family of dissipating moments}

For $a\in\R$, set
\[
 J_a(t)=\int_{\R}e^{ax}u(x,t)\,dx,
 \qquad
 \mathcal E_a(t)=e^{-(1+Da^2)t}J_a(t).
\]

\begin{theorem}[Exponential-moment dissipation]\label{thm:dissipation}
For every $a$ for which the preceding quantities are finite,
\begin{equation}
 \mathcal E_a'(t)
 =-\frac{e^{-(1+Da^2)t}}2
 \iint_{\R^2}K(x-y)\bigl(e^{ax}+e^{ay}\bigr)
 u(x,t)u(y,t)\,dx\,dy.
\end{equation}
In particular, $\mathcal E_a$ is nonincreasing.  It is strictly decreasing
whenever the double integral in \emph{(3)} is positive; this holds, for
example, when $u(\cdot,t)>0$ and $K$ is nonzero and nonnegative.
\end{theorem}

\begin{proof}
Twice integrating the diffusion term by parts gives
\begin{equation}
 J_a'=(1+Da^2)J_a-
 \int_{\R}e^{ax}u(x)(K*u)(x)\,dx.
\end{equation}
Evenness of $K$ permits exchange of $x$ and $y$, so
\begin{align*}
 \int e^{ax}u(x)(K*u)(x)\,dx
 &=\iint K(x-y)e^{ax}u(x)u(y)\,dx\,dy\\
 &=\frac12\iint K(x-y)(e^{ax}+e^{ay})u(x)u(y)\,dx\,dy.
\end{align*}
Differentiating the exponential normalization in $\mathcal E_a$ proves
\emph{(3)}.  Its integrand is nonnegative under the stated assumptions.
\end{proof}

For the top-hat kernel, (3) becomes
\begin{equation}
 \mathcal E_a'(t)
 =-\frac{e^{-(1+Da^2)t}}4
 \iint_{|x-y|\leq1}(e^{ax}+e^{ay})u(x,t)u(y,t)\,dx\,dy.
\end{equation}
At the critical tilts $a=\pm D^{-1/2}$ the normalizing exponent is $2$.

\section{Translation speed and competition delay}

Suppose that, for some $T>0$ and $\lambda>0$,
\begin{equation}
 u(x,t_0+T)=u(x-\lambda,t_0).
\end{equation}
Then $J_a(t_0+T)=e^{a\lambda}J_a(t_0)$.  Define the tilted
competition average
\begin{equation}
 \langle K*u\rangle_{a,t}
 =\frac{\int e^{ax}u(x,t)(K*u)(x,t)\,dx}
        {\int e^{ax}u(x,t)\,dx}.
\end{equation}

\begin{corollary}[Exact translation-delay identity]\label{cor:delay}
If $a>0$, $J_a$ is finite and nonzero on the interval, and \emph{(6)} holds,
then
\begin{equation}
 \frac{\lambda}{T}
 =Da+\frac1a-\frac1{aT}
   \int_{t_0}^{t_0+T}\langle K*u\rangle_{a,t}\,dt.
\end{equation}
Consequently,
\begin{equation}
 \frac{\lambda}{T}<Da+\frac1a
\end{equation}
whenever the competition integral is positive.  At $a=D^{-1/2}$,
\begin{equation}
 \frac{\lambda}{\sqrt D}
 =2T-\int_{t_0}^{t_0+T}\langle K*u\rangle_{D^{-1/2},t}\,dt,
 \qquad \frac{\lambda}{T}<2\sqrt D.
\end{equation}
\end{corollary}

\begin{proof}
Divide (4) by $J_a$, integrate in time, and use the translated endpoint
relation.  Strict positivity gives (9).  Substitution of $a=D^{-1/2}$ gives
(10).
\end{proof}

The first two terms on the right of (8) are the linear dispersion curve
$Da+a^{-1}$; its minimum over $a>0$ is the pulled speed $2\sqrt D$.  The last
term is an exact nonlinear phase delay, not an error estimate.

\section{What the obstruction says about patterned wakes}

Equation (10) rules out the following closure: a nonzero finite-critical-
moment state advances at pulled speed and, after each shedding time, becomes
an exact translate of itself.  The obstruction is not that a patterned wake
cannot form.  Rather, the complete state of a pulled invasion cannot be one
localized repeating packet.  At least one of the following must occur:
\begin{enumerate}
  \item the critical leading tail, whose critical exponential moment is
        typically divergent, is treated separately from the localized wake;
  \item the cycle state retains the time-integrated tilted competition trace
        in (10), or an equivalent memory variable;
  \item the shedding process is not exactly recurrent up to translation.
\end{enumerate}

This distinction matters because an invasion front can move at pulled speed
while leaving nearly stationary humps behind it.  Such a solution is not a
counterexample to Corollary~\ref{cor:delay}: its leading edge and wake are not
a single finite-moment translate.

\section{Open quantitative step}

The identity isolates a concrete next problem.  After decomposing a solution
into a critical leading component and a localized hump/wake component, one
must evaluate or bound
\[
 \int_{t_0}^{t_0+T}\langle K*u\rangle_{D^{-1/2},t}\,dt.
\]
If this integral becomes a universal functional of one established hump and
its incoming leading-edge amplitude, it could close a finite-dimensional
formation map.  If two nonlinear histories with the same visible hump retain
different values of this integral, then a one-hump Markov description is
impossible.  Neither alternative is proved here.

\section{Status and review questions}

Theorem~\ref{thm:dissipation} and Corollary~\ref{cor:delay} are
\textbf{PROVED} under their stated regularity and finiteness assumptions.
The interpretation as a necessary state decomposition is a direct logical
consequence of the corollary.  No claim of wavelength selection, asymptotic
periodicity, or biological prediction is made.

The author especially invites review of three points:
\begin{enumerate}
  \item Is identity (3), or its critical-tilt consequence (10), already
        explicit in the nonlocal Fisher--KPP literature?
  \item Can the regularity assumptions be weakened to a natural
        measure-valued or one-sided weighted class?
  \item Can the competition-delay integral be connected rigorously to a
        leading-edge/wake decomposition for top-hat kernels?
\end{enumerate}

\begin{thebibliography}{9}
\bibitem{HamelRyzhik2014}
F.~Hamel and L.~Ryzhik,
\newblock On the nonlocal Fisher--KPP equation: steady states, spreading
speed and global bounds,
\newblock \emph{Nonlinearity} \textbf{27} (2014), 2735--2753.
\newblock \url{https://doi.org/10.1088/0951-7715/27/11/2735}.

\bibitem{NeedhamEtAl2025}
D.~J. Needham, J.~Billingham, N.~M. Ladas, and J.~Meyer,
\newblock The evolution problem for the 1D nonlocal Fisher--KPP equation
with a top hat kernel. Part 1: The Cauchy problem on the real line,
\newblock \emph{European Journal of Applied Mathematics} \textbf{36}
(2025), 775--810.
\newblock \url{https://doi.org/10.1017/S0956792524000688}.
\end{thebibliography}

\end{document}
